\lhead{3. Knuth-Morris-Pratt (KMP) Algorithm}

\begin{enumerate}[label=\textbf{\Alph*.}]

\item \textbf{What is the KMP algorithm? State its main advantage over Brute Force
matching}

The Knuth-Morris-Pratt (KMP) Algorithm is an efficient pattern matching technique
that searches for a pattern within a text using previously matched information.

Instead of restarting comparisons after every mismatch, KMP intelligently decides
where to continue searching.

\textbf{Main Advantage}

It avoids repeated comparisons and performs pattern matching in $O(n + m)$ time,
making it much faster than Brute Force for large inputs.

\item \textbf{Explain the Prefix Function ($\Pi$ table) in KMP.}

The Prefix Function, also called the {\bfseries $\Pi$ Table or LPS (Longest
Prefix Suffix) Table}, stores the length of the longest proper prefix that is
also a suffix for each position in the pattern.

This information helps the algorithm decide how far the pattern should be shifted
when a mismatch occurs.

As a result, already matched characters do not need to be compared again.

\item \textbf{Describe how KMP avoids unnecessary comparisons during pattern searching.}

When a mismatch occurs, KMP does not move back to the beginning of the pattern.
Instead, it uses the values stored in the Prefix Table to determine the next comparison.
This saves time because characters that have already matched are not checked repeatedly.
Therefore, the searching process becomes much more efficient.

\item \textbf{Construct the Prefix Table for the following pattern and explain the process:\\
Pattern: ABABCABAB}

\begin{tabular}{|c|c|c|}
\hline
Position & Character & Value\\
\hline
0 & A & 0\\
1 & B & 0\\
2 & A & 1\\
3 & B & 2\\
4 & C & 0\\
5 & A & 1\\
6 & B & 2\\
7 & A & 3\\
8 & B & 4\\
\hline
\end{tabular}

\textbf{Final Prefix Table}: 0 0 1 2 0 1 2 3 4

\end{enumerate}

\begin{center}
\vspace{1cm}
$\ast \ast \ast$
\end{center}
\newpage
